% GENERATED FILE. Edit canonical lessons, then run npm run generate:books.
% canonical-source-hash: fnv1a64:7417333baa18070e
% canonical-lessons: TE-C288-aparesanu, TE-C288-ajirnam, TE-C288-malabaddhakam, TE-C288-madhumeham, TE-C288-raktapotu

\chapter{Operation, Indigestion, Constipation, Diabetes, Blood pressure}
\label{ch:te-operation-indigestion-constipation-diabetes-blood-pressure}
\hlchaptermodality{288}

\begin{chapteropening}
\textbf{By the end of this chapter:} \emph{I can say an operation, indigestion, constipation, diabetes and blood pressure.}

Complete the last lesson of chapter 288: I can say an operation, indigestion, constipation, diabetes and blood pressure.
\end{chapteropening}

\section[\texorpdfstring{\te{ఆపరేషను} (āparēṣanu)}{ఆపరేషను (āparēṣanu)}]{\te{ఆపరేషను} (āparēṣanu) — an operation, surgery}

\label{lesson:TE-C288-aparesanu}



\begin{quote}
Before the new one: say the Telugu for insurance, then the Telugu for an emergency.
\end{quote}

\subsection*{You'll want to know: \te{ఆపరేషను}}
\textbf{\te{ఆపరేషను}} — \emph{āparēṣanu} — \textquotedblleft{}an operation, surgery\textquotedblright{}.

\begin{grammarlens}[title={What you've built}]
One more word for this chapter.
\end{grammarlens}

\subsection*{Your turn}
\begin{itemize}
\raggedright
  \item \emph{Say it:} \emph{āparēṣanu}
  \item \emph{Say it:} \emph{āparēṣanu}, once more
  \item \emph{Say it:} say \emph{āparēṣanu}
  \item \emph{Recall:} say the Telugu for insurance, then the Telugu for an emergency, then say \emph{āparēṣanu} again
\end{itemize}

\begin{tcolorbox}[breakable,colback=teal!4,colframe=teal!35!black,title={Before you move on}]
What does \te{ఆపరేషను} mean? (\textquotedblleft{}An operation, surgery\textquotedblright{}.) Say it once more.
\end{tcolorbox}
\section[\texorpdfstring{\te{అజీర్ణం} (ajīrṇaṁ)}{అజీర్ణం (ajīrṇaṁ)}]{\te{అజీర్ణం} (ajīrṇaṁ) — indigestion}

\label{lesson:TE-C288-ajirnam}



\begin{quote}
Before the new one: say the Telugu for an emergency, then the Telugu for an operation.
\end{quote}

\subsection*{You'll want to know: \te{అజీర్ణం}}
\textbf{\te{అజీర్ణం}} — \emph{ajīrṇaṁ} — \textquotedblleft{}indigestion\textquotedblright{}.

\begin{grammarlens}[title={What you've built}]
One more word for this chapter.
\end{grammarlens}

\subsection*{Your turn}
\begin{itemize}
\raggedright
  \item \emph{Say it:} \emph{ajīrṇaṁ}
  \item \emph{Say it:} \emph{ajīrṇaṁ}, once more
  \item \emph{Say it:} say \emph{ajīrṇaṁ}
  \item \emph{Recall:} say the Telugu for an emergency, then the Telugu for an operation, then say \emph{ajīrṇaṁ} again
\end{itemize}

\begin{tcolorbox}[breakable,colback=teal!4,colframe=teal!35!black,title={Before you move on}]
What does \te{అజీర్ణం} mean? (\textquotedblleft{}Indigestion\textquotedblright{}.) Say it once more.
\end{tcolorbox}
\section[\texorpdfstring{\te{మలబద్ధకం} (malabaddhakaṁ)}{మలబద్ధకం (malabaddhakaṁ)}]{\te{మలబద్ధకం} (malabaddhakaṁ) — constipation}

\label{lesson:TE-C288-malabaddhakam}



\begin{quote}
Before the new one: say the Telugu for an operation, then the Telugu for indigestion.
\end{quote}

\subsection*{You'll want to know: \te{మలబద్ధకం}}
\textbf{\te{మలబద్ధకం}} — \emph{malabaddhakaṁ} — \textquotedblleft{}constipation\textquotedblright{}.

\begin{grammarlens}[title={What you've built}]
One more word for this chapter.
\end{grammarlens}

\subsection*{Your turn}
\begin{itemize}
\raggedright
  \item \emph{Say it:} \emph{malabaddhakaṁ}
  \item \emph{Say it:} \emph{malabaddhakaṁ}, once more
  \item \emph{Say it:} say \emph{malabaddhakaṁ}
  \item \emph{Recall:} say the Telugu for an operation, then the Telugu for indigestion, then say \emph{malabaddhakaṁ} again
\end{itemize}

\begin{tcolorbox}[breakable,colback=teal!4,colframe=teal!35!black,title={Before you move on}]
What does \te{మలబద్ధకం} mean? (\textquotedblleft{}Constipation\textquotedblright{}.) Say it once more.
\end{tcolorbox}
\section[\texorpdfstring{\te{మధుమేహం} (madhumēhaṁ)}{మధుమేహం (madhumēhaṁ)}]{\te{మధుమేహం} (madhumēhaṁ) — diabetes}

\label{lesson:TE-C288-madhumeham}



\begin{quote}
Before the new one: say the Telugu for indigestion, then the Telugu for constipation.
\end{quote}

\subsection*{You'll want to know: \te{మధుమేహం}}
\textbf{\te{మధుమేహం}} — \emph{madhumēhaṁ} — \textquotedblleft{}diabetes\textquotedblright{}.

\begin{grammarlens}[title={What you've built}]
One more word for this chapter.
\end{grammarlens}

\subsection*{Your turn}
\begin{itemize}
\raggedright
  \item \emph{Say it:} \emph{madhumēhaṁ}
  \item \emph{Say it:} \emph{madhumēhaṁ}, once more
  \item \emph{Say it:} say \emph{madhumēhaṁ}
  \item \emph{Recall:} say the Telugu for indigestion, then the Telugu for constipation, then say \emph{madhumēhaṁ} again
\end{itemize}

\begin{tcolorbox}[breakable,colback=teal!4,colframe=teal!35!black,title={Before you move on}]
What does \te{మధుమేహం} mean? (\textquotedblleft{}Diabetes\textquotedblright{}.) Say it once more.
\end{tcolorbox}
\section[\texorpdfstring{\te{రక్తపోటు} (raktapōṭu)}{రక్తపోటు (raktapōṭu)}]{\te{రక్తపోటు} (raktapōṭu) — blood pressure}

\label{lesson:TE-C288-raktapotu}



\begin{quote}
Before the new one: say the Telugu for constipation, then the Telugu for diabetes.
\end{quote}

\subsection*{You'll want to know: \te{రక్తపోటు}}
\textbf{\te{రక్తపోటు}} — \emph{raktapōṭu} — \textquotedblleft{}blood pressure\textquotedblright{}.

\begin{grammarlens}[title={What you've built}]
One more word for this chapter.
\end{grammarlens}

\subsection*{Your turn}
\begin{itemize}
\raggedright
  \item \emph{Say it:} \emph{raktapōṭu}
  \item \emph{Say it:} \emph{raktapōṭu}, once more
  \item \emph{Say it:} say \emph{raktapōṭu}
  \item \emph{Recall:} say the Telugu for constipation, then the Telugu for diabetes, then say \emph{raktapōṭu} again
\end{itemize}

\begin{tcolorbox}[breakable,colback=teal!4,colframe=teal!35!black,title={Before you move on}]
What does \te{రక్తపోటు} mean? (\textquotedblleft{}Blood pressure\textquotedblright{}.) Say it once more.
\end{tcolorbox}
